% BHU PYQ -- Transcribed & Corrected
% Paper: PHYSE-11 / PHYSE11: Mathematical Problem Solving Skills (Mid-Semester)
% Exam: B.Sc. (Hons.) Semester I Mid-Semester Examination, 2024-25 (NEP)
% Category: Skill Enhancement Course (SEC)
% Department: Physics

\documentclass[11pt,a4paper]{article}
\usepackage[a4paper,top=2.5cm,bottom=2.5cm,left=2.8cm,right=2.8cm,headheight=14pt]{geometry}
\usepackage{amsmath,amssymb}
\usepackage[T1]{fontenc}
\usepackage[utf8]{inputenc}
\usepackage{lmodern,microtype,fancyhdr,enumitem,hyperref,lastpage,parskip}

\hypersetup{
    colorlinks=true,
    urlcolor=black,
    linkcolor=black,
    pdftitle={PHYSE11: Mathematical Problem Solving Skills (Mid-Sem) -- BHU Sem I 2024-25 (NEP SEC)}
}

\pagestyle{fancy}
\fancyhf{}
\renewcommand{\headrulewidth}{0.4pt}
\renewcommand{\footrulewidth}{0.4pt}
\fancyhead[L]{\small\textit{Banaras Hindu University}}
\fancyhead[R]{\small\textit{PHYSE-11: Math Problem Solving Skills (Mid-Sem 2024-25)}}
\fancyfoot[C]{\small Page \thepage\ of \pageref{LastPage}}

\newcommand{\pts}[1]{\hfill\textit{[#1]}}

\begin{document}

\begin{center}
  {\large\bfseries BANARAS HINDU UNIVERSITY}\\[2pt]
  {\normalsize B.Sc. (Hons.) --- Semester I Mid-Semester Examination, 2024-25 (NEP)}\\[6pt]
  \rule{0.8\linewidth}{0.5pt}\\[6pt]
  {\large\bfseries PHYSICS (SEC)}\\[4pt]
  {\large\bfseries Paper Code: PHYSE-11 : Mathematical Problem Solving Skills (Mid-Sem)}\\[4pt]
  {\normalsize Time: 1 Hour\hfill Maximum Marks: 20}
\end{center}

\rule{\linewidth}{0.4pt}

\smallskip
\noindent\textit{Instructions: Answer all questions. Each question carries equal marks (5 marks each).}

\bigskip

\noindent\textbf{Question 1.} Use Taylor theorem to evaluate $\sqrt{1.5}$. Please limit yourself to using only the first 4 terms of the Taylor expansion.\pts{5}

\medskip\rule{\linewidth}{0.2pt}\medskip

\noindent\textbf{Question 2.} Draw the graph of function $f(x) = x^2 - 1$. Find the slope of the tangent line at the points $(-1,0)$, $(0,-1)$, and $(2,3)$.\pts{5}

\medskip\rule{\linewidth}{0.2pt}\medskip

\noindent\textbf{Question 3.} What is the difference between ordinary and partial differential equations? Explain with examples.\pts{5}

\medskip\rule{\linewidth}{0.2pt}\medskip

\noindent\textbf{Question 4.} Solve the following differential equation by using the integrating factor:
\[
  x \frac{dy}{dx} + 2y = 10x^2
\]\pts{5}

\vfill
\rule{\linewidth}{0.4pt}
\begin{center}\small
\href{https://bhu-exam.vercel.app/}{Click here to visit our website}\\[4pt]
\href{https://me-aryan.vercel.app/}{Click here to visit our contributor}
\end{center}

\end{document}
